% Suggested minimal manuscript revision. DO NOT compile into main.tex automatically.
% Recommended location: replace/expand the paragraph surrounding Eq. (27),
% at the end of Section III-B.

\hladd{Let $\Delta_i=|\mathcal{N}_i|$ denote the communication degree of
interceptor $i$. More explicitly, merging the Top-$L_{\max}$ records received
from its neighbors requires
$O\!\left(N\Delta_iL_{\max}\log(\Delta_iL_{\max})\right)$ operations and
$O(N\Delta_iL_{\max})$ received records per exchange. Because $L_{\max}$ is a
small fixed capacity, these terms reduce to $O(N\Delta_i)$ in
Eq.~\eqref{eq:complexity}. Consequently, for bounded communication degree and
fixed $N$, $P$, and $T$, the total per-round computation and communication of
the swarm grow linearly with the interceptor number $M$.}

\hladd{For a connected graph of diameter $D$, a winning bid requires at most
$D$ hops to propagate under a static assignment snapshot. If the maximum
single-hop delay is $d_{\max}$ communication rounds, the conservative CBBA
consensus bound becomes
$O\!\left(ND(d_{\max}+1)\right)$ exchange rounds, or
$O\!\left(ND(d_{\max}+1)T_{\mathrm{ex}}\right)$ in wall-clock time, where
$T_{\mathrm{ex}}$ is the communication-exchange period. Thus, bounded delay
does not alter the static Top-$L_{\max}$ fixed point when the graph remains
connected and messages are eventually delivered, but it increases the
consensus time. In a time-varying engagement, the allocation snapshot must be
updated on a time scale longer than the consensus interval to avoid decisions
based on stale bids.}

